\input eplain
\pdfpagewidth=148mm
\pdfpageheight=210mm
\hsize=128mm
\vsize=190mm

\pdfhorigin=0pt
\pdfvorigin=0pt
\hoffset=10mm
\voffset=10mm
\noindent news \leaders\hbox to 10pt {\hss.\hss}  \hfill \xref{node1}
\smallskip
\noindent 05 limits \leaders\hbox to 10pt {\hss.\hss}  \hfill \xref{node0}
\smallskip
\vfill \break
{\medskip\noindent\bf\font\titlefont=cmss17 at 20pt\titlefont Timeline}\medskip
{\noindent \bf 2026-01-28:} news, 05 limits\smallskip
{\noindent \bf 2026-01-29:} 05 limits\smallskip
{\noindent \bf 2026-01-30:} 05 limits\smallskip
{\noindent \bf 2026-02-03:} 05 limits\smallskip
{\noindent \bf 2026-02-04:} 05 limits\smallskip
{\noindent \bf 2026-02-11:} 05 limits\smallskip
\vfill \break
\xrdef{node0}
{\medskip\noindent\bf\font\titlefont=cmss17 at 20pt\titlefont 05 limits}\medskip
{\medskip\noindent\bf\font\titlefont=cmitt12 at 17pt\titlefont 2026-01-28}\medskip
{\smallskip\noindent\bf\font\titlefont=cmr12 at 13pt\titlefont $\bullet$ 09:00: some initial spivak reading}\smallskip
Since I'm reading about limits in my calc class, I thought I'd crack this one open and take a stab at it. It provides a completely different mental model for thinking about limits. Definitely needed to read things carefully to be able to absorb what was on the page (I really only read 2 pages).
\par

{\medskip\noindent\bf\font\titlefont=cmitt12 at 17pt\titlefont 2026-01-29}\medskip
{\smallskip\noindent\bf\font\titlefont=cmr12 at 13pt\titlefont $\bullet$ 09:53: limits}\smallskip
Focusing on $x\sin(1/x)$.
\par

{\smallskip\noindent\bf\font\titlefont=cmr12 at 13pt\titlefont $\bullet$ 10:12: limits}\smallskip
After $x\sin(1/x)$, examine $x^2\sin(1/x)$ as it approaches near 0. Begin to consider function $\sqrt(x)\sin(1/x)$.
\par

{\medskip\noindent\bf\font\titlefont=cmitt12 at 17pt\titlefont 2026-01-30}\medskip
{\smallskip\noindent\bf\font\titlefont=cmr12 at 13pt\titlefont $\bullet$ 08:50: limits}\smallskip
Examining $\sin(1/x)$. False that it approaches zero near zero, with proof (?). Then it went to show that it doesn't approach any number near 0. The thing that's tripping me up is the statement about intervals: for any interval A containing 0, there is some number containing $x = 1 / (90 + 360n)$ which is in this interval, for which $f(x) = 1$.
\par

{\medskip\noindent\bf\font\titlefont=cmitt12 at 17pt\titlefont 2026-02-03}\medskip
{\smallskip\noindent\bf\font\titlefont=cmr12 at 13pt\titlefont $\bullet$ 07:25: spivak reading}\smallskip
Function that checks if number is rational or not, analyzing the limit.
\par

{\medskip\noindent\bf\font\titlefont=cmitt12 at 17pt\titlefont 2026-02-04}\medskip
{\smallskip\noindent\bf\font\titlefont=cmr12 at 13pt\titlefont $\bullet$ 09:51: spivak reading}\smallskip
An "amusing variation" on function to determine if function is rational or not. Simpler functions. Constants. A sign checker function.
\par

{\medskip\noindent\bf\font\titlefont=cmitt12 at 17pt\titlefont 2026-02-11}\medskip
{\smallskip\noindent\bf\font\titlefont=cmr12 at 13pt\titlefont $\bullet$ 10:06: spivak reading}\smallskip
At this point in the reading, the provision definition of a limit is revisited, as it is not precise enough to be used in proofs. The proof is revised by fixing one ambiguity at a time.
\par

\xrdef{node1}
{\medskip\noindent\bf\font\titlefont=cmss17 at 20pt\titlefont news}\medskip
{\medskip\noindent\bf\font\titlefont=cmitt12 at 17pt\titlefont 2026-01-28}\medskip
{\smallskip\noindent\bf\font\titlefont=cmr12 at 13pt\titlefont $\bullet$ 18:59: logging, seting up some new things related to my calculus class}\smallskip
\bye
